\documentclass[UTF8,12pt,a4paper,fontset=fandol]{ctexart}

% 支持从项目根目录、深度学习目录或当前笔记目录编译。
\makeatletter
\def\input@path{{../}{../../}}
\makeatother
\usepackage{researchnotes}

\title{类型安全决策模型与Jev}
\author{}
\date{}

\begin{document}
\maketitle

\section{从文本生成到类型化决策}
\label{sec:jev-motivation}

退款分类、工单路由等任务需要程序可直接使用的判断。若仅用提示词要求大语言模型
（large language model, LLM）输出 JSON，仍可能出现额外说明、括号缺失或非法枚举值。
\keyterm{自回归生成}（autoregressive generation）逐个预测词元（token），
词元并不等于单个汉字；提示词中的格式要求也不自动成为生成约束。
不过，\keyterm{约束解码}（constrained decoding）可以在生成时排除不符合规则的词元，
因此不能把格式错误说成所有 LLM 结构化输出的必然缺陷，参见
\href{https://openai.com/index/introducing-structured-outputs-in-the-api/}{结构化输出技术说明}。

TypeSafe 将 Jev 定位为\keyterm{类型安全决策模型}：输入业务状态与问题，
直接返回预先规定类型的值及概率，不自由生成自然语言答案。
\href{https://typesafe.ai/blog/introducing-system-one-models-and-jev}{官方发布说明}
强调并行输出与面向校准决策的训练方法；这支持其接口定位，但不足以断言未公开的解码头结构。
API 中仍可出现 JSON 字段名和预设选项字符串，它们不等于模型自由撰写的文本。

\section{三种输出原语及其含义}
\label{sec:jev-primitives}

\keyterm{原语}（primitive）是可由程序组合的基本判断单元。下表概括三类原语，
具体限制以官方文档为准。

\begin{center}
\begin{tabularx}{\textwidth}{@{}lYY@{}}
\toprule
原语 & 判断目标 & 输出及边界 \\
\midrule
\href{https://docs.typesafe.ai/primitives/noul}{Noul}
& 命题是否为真 & 返回 $[0,1]$ 内的概率；无单独的置信度字段。 \\
\href{https://docs.typesafe.ai/primitives/choice}{Choice}
& 从给定类别中选择一项 & 最多 255 个选项；返回所选项、完整概率分布及置信度。 \\
\href{https://docs.typesafe.ai/primitives/score}{Score}
& 按有序量表评定程度 & 定义 2--10 个档位；返回分数、档位说明、概率分布及置信度。 \\
\bottomrule
\end{tabularx}
\end{center}

设 $x$ 为输入状态，$q$ 为包含判定标准的问题。
Noul 可记为 $p_\theta(Y=1\mid x,q)$，其中 $Y\in\{0,1\}$ 表示命题真假，
$\theta$ 表示模型参数。这是模型的条件概率估计，不能直接视为真实后验概率。
例如 $0.82$ 表示模型给“真”分配的概率，并非“满足条件的程度为 82\%”。

Choice 在有限集合 $\mathcal C=\{c_1,\ldots,c_K\}$ 上给出概率 $p_k$，
满足 $p_k\geq0$、$\sum_k p_k=1$，并返回概率最大的选项。
它实现\keyterm{闭集分类}（closed-set classification）；类别不完备时应设“其他”选项。
Score 的档位编号为 $0,\ldots,m-1$，其档位概率记为 $r_j$，满足
$r_j\geq0$、$\sum_j r_j=1$。返回分数为
\begin{equation}
  s=\sum_{j=0}^{m-1}j\,r_j\in[0,m-1].
  \label{eq:jev-score}
\end{equation}
所以档位离散，但分数不必为整数。例如档位概率为 $(0.1,0.3,0.6)$ 时，
$s=1.5$。这是档位编号的加权均值，并不自动具有物理量的等距含义。

\section{类型安全与业务正确性}
\label{sec:jev-reliability}

\keyterm{类型安全}（type safety）保证输出属于约定范围，不能保证选对类别、
读懂政策或输入资料真实。官方的 \texttt{confidence} 是由概率分布计算的集中程度指标，
不能直接当作“答案正确的概率”，参见
\href{https://docs.typesafe.ai/confidence}{置信度说明}。
\keyterm{概率校准}（probability calibration）要求在大量预测中，
获得相近预测概率的事件，其实际发生频率与该概率相近；它需要在目标业务数据上检验。

以退货分流为例，可让 Noul 判断用户是否仅因改变主意而退货，
让 Choice 区分质量问题、个人偏好和其他原因，让 Score 评估材料完整程度。
程序再结合订单期限等确定性规则决定处理路径，并将不确定案例转人工复核。
只依赖同一输入的问题可并行提交；依赖前一答案的问题需另行调用。
这样可以减少自由文本修复和重复生成的开销，但仍需评估误判、延迟和异常处理成本。

\section{如何阅读成本与准确率评测}
\label{sec:jev-evaluation}

截至 2026 年 9 月 29 日，\href{https://evals.typesafe.ai/}{官方工作流评测}
涵盖安全事件、智能体运行记录分析、发票处理和客户服务。
总览对四类工作流等权平均，显示 Jev 约为每例 0.0004 美元、67.8\% 准确率，
平均用时约 0.4 秒；原材料中的 0.0003 美元不宜作为固定单价。
同图 Terra 的成本约为 0.0304 美元、准确率为 67.9\%，
说明此配置下的成本优势接近两个数量级，而非对所有 GPT、Claude 调用都成立。

参考标签来自 GPT-6 Astra 与 Claude Fable 5.1 在高思考设置下的回答共识，
其余模型使用服务商默认推理设置。因此，该指标反映与参考共识的一致程度，
不能等同于独立人工真值准确率，也不能统称为与“关闭思考的轻量模型”比较。
若横轴为成本、纵轴为准确率，\keyterm{帕累托前沿}（Pareto frontier）
由不存在另一配置成本不更高且准确率不更低、并至少一项严格更优的点组成。
Jev 位于图中前沿的低成本端，并非准确率最高点；实际选型仍取决于业务误判代价。

\end{document}
